\documentclass[12pt]{article}
\usepackage[utf8]{inputenc}
\usepackage[T1]{fontenc}
\usepackage{hyperref}
\title{Yawalapit\'i Language Resources for K-12 Curriculum Development}
\author{Compiled by Claude Sonnet 5.5 (claude-sonnet-5-5) for Andr\'as Kornai}
\date{}
\begin{document}
\maketitle

\section{Introduction}
Yawalapit\'i (ISO 639-3: yaw; Glottocode yawa1261) is an Arawakan language of the Xingu Indigenous Park, Mato Grosso, Brazil. It belongs to the Xinguan Arawak group, close to Waur\'a and Mehinaku. Speaker counts in the sources found are very low and vary: Wikipedia gives 8 (2006), the Povos Ind\'igenas no Brasil site gives four or five, and the Endangered Languages Project mentions three full speakers. Searches found no dictionary, primer or school material.

\section{Primary Educational Resources}
No primers, textbooks or school programmes for Yawalapit\'i were found. The reference description says the language is not known to be taught in schools.

\section{Linguistic and Reference Materials}
\begin{itemize}
\item Wikipedia article on Yawalapit\'i, which covers phonology and morphology: \url{https://en.wikipedia.org/wiki/Yawalapit%C3%AD_language}
\item Glottolog entry for the language: \url{https://glottolog.org/resource/languoid/iso/yaw}
\item Endangered Languages Project entry, which cites a Portuguese-language documentation project and a 2007 description by Mily Crevels: \url{https://endangeredlanguages.com/lang/1639}
\end{itemize}

\section{Religious and Community Resources}
\begin{itemize}
\item Povos Ind\'igenas no Brasil (Instituto Socioambiental) page on the Yawalapit\'i people: \url{https://pib.socioambiental.org/en/povo/Yawalapiti}
\item Indios.org.br page on the Yawalapit\'i: \url{https://indios.org.br/en/povo/yawalapiti}
\end{itemize}

\section{Digital and Audio Resources}
No audio, video or digital resources were found in the searches.

\section{Curriculum Development Context}
The language is critically endangered, with only a handful of elderly speakers. A curriculum would be a revitalization effort based on documentation and work with the remaining speakers and the community, with support from linguists. The Crevels description and the documentation project cited by the Endangered Languages Project are the leads to follow. Almost nothing is available in the sources found.

\end{document}